\section{Related Work}
\label{sec:related_work}

\paragraph{DIF Methodology and Measurement Invariance.}
The Mantel-Haenszel procedure for DIF detection was introduced by \citet{holland1988mh} and remains the most widely used method in educational testing due to its simplicity and well-characterized statistical properties \citep{clauser1998dif}. Alternative approaches include logistic regression DIF \citep{swaminathan1990logistic}, SIBTEST \citep{shealy1993sibtest}, Lord's $\chi^2$ test \citep{lord1980irt}, and the Raju area method \citep{raju1988area}; \citet{millsap1993dif}, \citet{zumbo1999handbook}, and \citet{penfield2007dif} provide comprehensive reviews. DIF analysis is a specific instance of measurement invariance testing \citep{meredith1993invariance, vandenberg2000invariance, millsap2011invariance}, which examines whether a measurement instrument functions equivalently across populations. MIMIC models offer a latent-variable approach to the same problem \citep{muthen1985mimic, woods2009mimic}, and \citet{teresi2007dif} reviews DIF methods across applied domains. Cross-subject matching---computing the stratification variable from domains other than the one being tested---adapts purified subtest matching from the DIF literature \citep{nandakumar1993dif,donoghue1993mh} to a new challenge: contamination control at scale, where the matching variable itself may be compromised by the phenomenon under test. We retain the standard MH statistic and develop LLM-specific adaptations around it: cross-subject matching for contamination control (with a constructive $\Delta\mathrm{FPR} = 0$ guarantee), DEFF correction for non-independence, and semi-synthetic validation absent ground-truth labels. The 5,227-examinee scale exceeds typical DIF applications~\citep{belzak2020smallsample}. The impact-versus-DIF confound is addressed through matched-ability quintile analysis, IRT $\theta$-conditioning (as sensitivity check), and a random-split control (\Cref{sec:temporal}).

\paragraph{IRT in LLM Evaluation.}
IRT adoption in LLM evaluation has accelerated: \citet{lalor2019irt} introduced IRT to NLP, \citet{vania2021irt} assessed test set discriminability, and \citet{rodriguez2021evaluation} demonstrated that not all evaluation examples are equally informative. Recent extensions include PSN-IRT \citep{zhou2026psnirt} for benchmark quality at scale, TinyBenchmarks \citep{maiapolo2024tinybenchmarks} for IRT-based item selection, scalability coefficients for detecting bad benchmark items \citep{hardy2026scalability}, IRT-based evaluation quality~\citep{kraus2024irt}, ATLAS \citep{li2025atlas} for computerized adaptive testing of LLMs, JE-IRT \citep{yao2025jeirt} for geometric ability embedding, continuous-score adaptive evaluation~\citep{balkir2026confident}, and competency-focused IRT evaluation in medical domains \citep{luo2025medirt}. The closest concurrent work is Growing Pains \citep{habba2026growingpains}, which uses IRT calibration to establish anchor items for cross-temporal score equating via fixed-parameter calibration \citep[FPC;][]{kim2019fpc}. While Growing Pains uses parametric IRT (2PL/3PL fixed-parameter calibration), MH-DIF is nonparametric and robust to model misspecification---an advantage given poor 2PL fit for MMLU (RMSEA $> 0.08$; \Cref{app:power_sensitivity}). Growing Pains prescribes how to equate scores after instability is detected; our work diagnoses \emph{which} items exhibit instability and quantifies severity via ETS classification, providing pre-screened anchor candidates for downstream equating. Their anchor-item selection can directly consume our stable-item set as pre-screened candidates; our finding that 31\% of items exhibit temporal DIF identifies which items should \emph{not} serve as anchors. To our knowledge, no prior work applies DIF to benchmark temporal comparability---contamination detection and psychometric evaluation optimization have developed as separate threads.

\paragraph{Contamination Detection.}
A growing family of methods detects benchmark contamination through model behavior or training data analysis. Membership inference approaches estimate whether a model has seen specific items during training: Min-K\% \citep{shi2024mink} and Min-K\%++ \citep{zhang2024minkpp} use token-level log-likelihood statistics, while Time Travel \citep{golchin2024timetravel} probes models with temporal cues and \citet{oren2024proving} prove contamination through exchangeability tests. Performance-based detectors identify contamination through statistical anomalies \citep{dekoninck2024constat,liang2026dvd,tao2026selfcritique,deng2024investigating} in-context learning effects~\citep{zawalski2026codec}, or inference-time decontamination~\citep{chai2026benchmarksleak}. \citet{balloccu2024leak} provide a broader taxonomy of leakage detection strategies. These methods share a common validation practice: they verify detection accuracy against web-overlap or n-gram labels \citep{li2024contamination} that classify items based on verbatim matches in publicly accessible corpora. Such labels measure \emph{global exposure} (whether an item exists in training data) rather than \emph{differential item functioning} (whether an item behaves differently across model cohorts). This global-vs-differential distinction motivates our comparison with web-overlap labels (\Cref{sec:groundtruth}). Work on LLM memorization \citep{carlini2022quantifying, tirumala2022memorization} establishes that memorization depends on data frequency and model scale, reinforcing the distinction between global exposure and differential behavior. \citet{wang2026fragility} show that existing contamination detectors are fragile under reinforcement learning, further motivating behavioral approaches grounded in psychometric theory.

\paragraph{Benchmark Quality and Design.}
Meta-evaluation of static benchmark quality~\citep{zhu2026bhi,qian2026benchmark2} assesses quality at a single snapshot; our temporal DIF analysis diagnoses \emph{when} measurement properties shift. MMLU-CF~\citep{zhao2025mmlucf} reconstructs contamination-free items and MMLU-Pro~\citep{wang2024mmlupro} extends item difficulty; DIF monitoring identifies which existing items remain measurement-invariant without replacement. \citet{maeda2025difllm} use LLMs to predict DIF in human assessments---the reverse direction; we apply DIF to LLM benchmarks with models as examinees. Dynamic benchmarks~\citep{kiela2021dynabench,li2024latesteval,white2025livebench} and self-evolving frameworks~\citep{ying2024automating,xu2025benchmark} mitigate contamination by construction but lack psychometric equating. Broader surveys~\citep{chang2024survey,burnell2023rethink,alzahrani2024benchmarks} identify benchmark saturation as a systemic challenge that DIF analysis directly addresses.
